\chapter{MATERIALS AND METHODS}

\section{Study design and data source}

This was a cross-sectional secondary analysis of the Viet Nam Multiple Indicator Cluster Survey 2020--2021 (MICS6), conducted by the General Statistics Office of Viet Nam with technical support from UNICEF. MICS6 is a nationally representative household survey using a stratified, two-stage cluster sampling design based on the 2019 Viet Nam Population and Housing Census. The sample was stratified by urban and rural areas, socioeconomic region, major cities, and ethnic-minority composition of enumeration areas. In the first stage, 700 enumeration areas were selected with probability proportional to size; in the second stage, 20 households were systematically selected within each enumeration area, yielding 14,000 selected households. The women's questionnaire was administered to eligible women aged 15--49 years in selected households; the full women's module interviewed 10,770 women before this thesis restricted the sample to women aged 15--29 years \cite{unicef2022}.

The analysis used the women's individual dataset in international Stata format. All analyses accounted for the complex survey design, including sampling weights, primary sampling units, and strata.

\section{Study population}

The target population for this thesis was Vietnamese women aged 15--29 years. This age range describes older adolescents and young adult women before broad public financing of HPV vaccination; it should not be interpreted as estimating coverage among the standard primary target group of girls aged 9--14 years. The analytic sample was defined by the following criteria:

\begin{itemize}
    \item Women in MICS age groups WAGE 1--3, corresponding to ages 15--19, 20--24, and 25--29 years.
    \item Non-missing and positive women's sampling weight (\texttt{wmweight}).
    \item Non-missing primary sampling unit (\texttt{PSU}) and stratum (\texttt{stratum}).
\end{itemize}

After restricting the dataset to women aged 15--29 years and excluding records without valid survey design information, the final analytic sample included \textbf{4,102} women. No additional sample size calculation was performed because the study used all eligible observations in the nationally representative dataset.

\begin{figure}[H]
\centering
\begin{tikzpicture}[
    >=Latex,
    flowStep/.style={draw, rounded corners, align=center, text width=9.2cm, minimum height=0.95cm, fill=gray!8},
    result/.style={draw, rounded corners, align=center, text width=9.2cm, minimum height=1.55cm, fill=green!8},
    arrow/.style={->, thick, shorten >=2pt, shorten <=2pt}
]
\node[flowStep] (all) at (0,0) {Women's module: all women aged 15--49 years};
\node[flowStep] (age) at (0,-1.60) {Restrict to age 15--29 years\\MICS age groups WAGE 1--3};
\node[flowStep] (design) at (0,-3.25) {Exclude records with missing/invalid\\sampling weight, PSU, or stratum};
\node[result] (final) at (0,-5.95) {\textbf{Final analytic sample: N = 4,102 women aged 15--29 years}\\Heard of HPV vaccination: 62.4\%\\Received HPV vaccine: 7.5\%\\Aware but unvaccinated among aware women: 88.0\% (aware N = 2,133)};
\draw[arrow] (all) -- (age);
\draw[arrow] (age) -- (design);
\draw[arrow] (design) -- (final);
\end{tikzpicture}
\caption{Study design and participant selection. The analytic sample included 4,102 women aged 15--29 years from Viet Nam MICS6 with valid survey design variables.}
\label{fig:study_flowchart}
\end{figure}

\section{Outcome variables}

Three primary binary outcomes and one sensitivity outcome were constructed:

\begin{enumerate}
    \item \textbf{HPV vaccine awareness} (\texttt{ccp5\_bin}): coded 1 if the respondent reported having heard of HPV vaccination and 0 otherwise.
    \item \textbf{Self-reported any-dose HPV vaccination status in the full analytic sample} (\texttt{ccp6\_bin}): coded 1 if the respondent reported having received HPV vaccine and 0 otherwise. Because the vaccination question was asked conditional on awareness in the MICS skip pattern, full-sample vaccination prevalence should be interpreted as self-reported any-dose uptake under the conservative assumption that women who had never heard of HPV vaccination had not received it. This outcome does not measure dose completion, vaccine product, timing, payment source, or provider type.
    \item \textbf{Awareness--vaccination gap} (\texttt{gap}): defined only among women who had heard of HPV vaccination. It was coded 1 for aware but unvaccinated and 0 for aware and vaccinated.
    \item \textbf{HPV vaccination among aware women} (\texttt{ccp6\_aware}): coded 1 for vaccinated and 0 for unvaccinated among women who had heard of HPV vaccination. This complementary outcome was analyzed explicitly to assess whether uptake patterns persisted after conditioning on awareness and to separate information barriers from post-awareness access barriers.
\end{enumerate}

\section{Explanatory variables and confounder measures}

Variables were selected based on the conceptual framework, prior literature, and availability in MICS6. The adjusted models were intended to estimate conditional associations rather than total causal effects. Because structural variables such as education, household wealth, ethnicity, and residence may lie on overlapping causal pathways, estimates should not be interpreted as causal effects of isolated variables. Table~\ref{tab:variables} lists all explanatory variables used in descriptive and bivariate analyses. ICT exposure was included descriptively and bivariately, but it was excluded from adjusted regression and decomposition because its weighted prevalence was 97.9\%, leaving sparse contrast and a high risk of unstable estimates.

\begin{table}[htbp]
\caption{Explanatory variables and operational definitions}
\centering
\small
\begin{tabular}{|p{3.3cm}|p{4.3cm}|p{5.6cm}|}
\hline
\textbf{Variable} & \textbf{Categories} & \textbf{Operational definition} \\
\hline
Age group & 15--19, 20--24, 25--29 & MICS age group variable WAGE \\
\hline
Residence & Urban, rural & Household area variable \\
\hline
Ethnicity & Kinh/Hoa, ethnic minority & MICS two-category ethnicity variable \\
\hline
Marital status & Currently married/in union, formerly married/in union, never married/in union & MICS marital status variable \\
\hline
Education & Primary or less, lower secondary, upper secondary or higher & Collapsed from MICS education level \\
\hline
Wealth quintile & Q1 poorest to Q5 richest & MICS household wealth quintile \\
\hline
Health insurance & Insured, uninsured & Recode of health insurance variable \\
\hline
Mass media exposure & No, yes & Any exposure to newspaper, radio, or television \\
\hline
ICT exposure & No, yes & Any reported computer, internet, or mobile phone exposure/use item in the media module \\
\hline
Gender-equitable attitude & No, yes & Yes if the respondent rejected all five wife-beating justifications \\
\hline
Sexual experience & No, yes & Ever had sexual intercourse \\
\hline
\end{tabular}
\label{tab:variables}
\end{table}

\section{Statistical analysis}

All statistical analyses were performed in StataNow 19.5 SE. The survey design was declared as:

\begin{flushleft}
\small\ttfamily
svyset PSU [pweight=wmweight], strata(stratum)\\
vce(linearized) singleunit(centered)
\end{flushleft}

\subsection{Descriptive analysis}

Survey-weighted percentages and 95\% confidence intervals were estimated for all outcomes and explanatory variables. These estimates used the \texttt{svy} prefix and therefore accounted for weights, clustering, and stratification.

\subsection{Bivariate analysis}

For each explanatory variable, outcome prevalence was estimated using survey-weighted proportions. Bivariate p-values were obtained from design-based Pearson chi-squared tests from \texttt{svy: tabulate, pearson}. These tests were descriptive, emphasized alongside effect sizes and 95\% confidence intervals, and were not adjusted for multiple comparisons. For the gap outcome, all bivariate analyses used \texttt{svy, subpop(ccp5\_bin)} so that inference was restricted to women who had heard of HPV vaccination while preserving the original survey design.

\subsection{Adjusted regression models}

Adjusted prevalence ratios were estimated using survey-weighted Poisson regression with a log link. This approach was chosen because odds ratios can overstate relative associations when outcomes are common, particularly for HPV vaccine awareness \cite{barros2003,zou2004}. The \texttt{svy} prefix used Taylor linearized survey variance estimation, accounting for stratification, clustering, and sampling weights; this provides design-based standard errors for the modified Poisson models. Models included age group, residence, ethnicity, marital status, education, wealth quintile, insurance, mass media exposure, gender-equitable attitude, and sexual experience. ICT exposure was excluded from adjusted models because of near-universal prevalence. Adjusted models used complete-case records for the modeled outcome and covariates; the main reduction in model sample size was due to missing gender-equitable attitude information.

For awareness and vaccination, models were fitted in the full analytic sample. For the awareness--vaccination gap and vaccination-among-aware sensitivity outcome, models used the aware subpopulation while preserving the survey design. Logistic regression models estimating odds ratios were conducted as sensitivity analyses and reported in the appendix. Adjusted predicted probabilities by wealth quintile were also estimated from survey-weighted logistic models to present absolute differences alongside prevalence ratios. Adjusted models used complete-case analysis; item missingness for model covariates and the distribution of complete-case exclusions are exported by the Stata pipeline as \texttt{tableA\_missingness.csv} and \texttt{tableA\_complete\_case\_profile.csv}.

\subsection{Concentration index and decomposition}

Socioeconomic ranking used the continuous household wealth score (\texttt{ses}). Wealth quintiles were used for grouped descriptive analysis and as dummy variables in adjusted and decomposition models. Because wealth quintiles were also derived from the wealth distribution used to rank inequality, their decomposition contributions should be interpreted as descriptive accounting of the wealth gradient rather than as independent actionable causes.

The standard concentration index was estimated for awareness, vaccination, and the awareness--vaccination gap. Because all outcomes were binary, the Erreygers-corrected concentration index was also estimated. Positive values indicate pro-rich concentration; negative values indicate pro-poor concentration.

Decomposition of the Erreygers index used a probit model to obtain average marginal effects for explanatory variables. Categorical covariates were entered as dummy variables, with reference categories matching the adjusted regression models: urban residence, age 25--29 years, currently married/in union, upper secondary or higher education, insured, Kinh/Hoa ethnicity, richest wealth quintile, no mass media exposure, no gender-equitable attitude, and no sexual experience. For each covariate dummy, its own survey-weighted standard concentration index was estimated using the continuous wealth score as the rank variable. The contribution of each covariate was calculated as:

\begin{equation}
\text{Contribution}_k = 4 \times AME_k \times \bar{x}_k \times CI_k,
\end{equation}

\noindent where $AME_k$ is the average marginal effect, $\bar{x}_k$ is the weighted mean of the covariate, and $CI_k$ is the covariate's standard concentration index. Percentage contribution was calculated by dividing each contribution by the total Erreygers index. Bootstrap confidence intervals for decomposition contributions were not estimated; therefore, decomposition results are presented as approximate explanatory accounting of the observed inequality pattern and not as causal attribution or statistically tested differences between contributors. A useful future sensitivity analysis would repeat decomposition after excluding wealth-quintile indicators to identify non-wealth covariates that remain aligned with the wealth gradient.

\section{Ethical considerations}

This study used de-identified secondary data from MICS6. The original survey obtained ethical approval and respondent consent according to UNICEF and national survey procedures \cite{unicef2022}. No direct contact with human participants occurred in this thesis. Any requirement for additional institutional review or exemption at Hanoi Medical University should be documented according to the university's thesis procedures before final submission.

\clearpage
